\subsection{Sprint 0}

A primeira sprint teve como objetivo principal realizar as seguintes atividades: promover a integração da equipe, estabelecendo uma base sólida de colaboração; participar ativamente do encontro com a stakeholder para elicitação de requisitos; definir as tecnologias mais adequadas para o projeto; modelar o banco de dados; desenvolver os protótipos das telas; e, finalmente, apresentar os resultados alcançados à stakeholder.

No início da sprint, a equipe concentrou seus esforços na integração, criando um ambiente de trabalho colaborativo e harmonioso. Em seguida, foram realizadas reuniões para discutir as ideias iniciais e planejar os próximos passos. Participei ativamente da formulação de perguntas, contribuindo com um documento colaborativo criado pelo time para o encontro com a stakeholder.

A reunião com a stakeholder foi extremamente esclarecedora, permitindo que a equipe definisse com mais clareza as arquiteturas e tecnologias a serem utilizadas no projeto. A partir disso, a equipe se organizou em dois grupos: um responsável pela prototipação do aplicativo para usuários e outro pelo dashboard de pesquisador. Utilizamos a plataforma Figma para criar o fluxo de telas e fiquei encarregada da prototipação da tela do dashboard, que desenvolvi e finalizei em colaboração com minha dupla.

Além disso, participei da reunião de modelagem do banco de dados, contribuindo com sugestões e levantando questões importantes. Durante essa fase, também auxiliei colegas com dificuldades em suas telas e configurei meu ambiente de desenvolvimento para iniciar as próximas etapas. Um dos maiores desafios identificados nesta sprint foi o desenvolvimento do player de áudio com funções de contagem de tempo, uma funcionalidade crucial para o projeto. Diante dessa necessidade, dediquei-me a pesquisar alternativas e encontrei uma biblioteca que poderia facilitar o desenvolvimento desse componente: ReactPlayer \cite{ReactPlayer}. Desenvolvi um protótipo que atendeu plenamente às especificações funcionais e de interface, eliminando a necessidade de upload de arquivos de áudio, o que otimiza tanto o tempo de desenvolvimento quanto os custos do projeto. Disponibilizei o código no repositório para que todos na equipe pudessem acessar e aprimorar o componente conforme necessário.

Entre os principais desafios enfrentados, destacaram-se a complexidade do envio de notificações personalizadas e a decisão sobre qual tecnologia adotar para o player de áudio. Além disso, foi identificado que o desenvolvimento de um aplicativo mobile traria dificuldades adicionais de integração, o que levou à decisão de substituir essa ideia por uma aplicação web responsiva. Essa mudança reduziu custos e facilitou o acesso dos usuários, eliminando a necessidade de download.

Ao longo dessa sprint, ficou claro que a discussão e a análise conjunta de diferentes abordagens são fundamentais para a tomada de decisões acertadas. Com a prototipação finalizada e as tecnologias definidas, o próximo passo será dar início à implementação do código, avançando com as funcionalidades necessárias para cumprir os requisitos do projeto.